\documentclass[letterpaper,12pt]{article}
%\usepackage{harvard}
%\usepackage{subfigure}
%\usepackage[active]{srcltx}
\usepackage{graphicx}
\usepackage{pdflscape}
\usepackage{amsmath}
\usepackage{amsthm}
\usepackage{lscape}
\usepackage{enumerate}
\usepackage{float}
\usepackage{grffile}
\usepackage{relsize}
\usepackage{tablefootnote}
\usepackage{footnote}
\usepackage{sectsty}
\usepackage{multirow}
\usepackage{xcolor}
\usepackage{caption}
\usepackage{color}
\usepackage[breaklinks,colorlinks]{hyperref}
\usepackage{natbib}
\usepackage{multirow}
\usepackage{bbm}
\usepackage[T1]{fontenc}
\usepackage[multiple]{footmisc}
\usepackage{appendix}
\usepackage{times}
\usepackage{relsize}
\usepackage{caption}
\usepackage{newtxtext,newtxmath}
%%%%%%%%%%%%%%%%%%%%%%%%%%%%%%%%%%%%%%%%%%%%%%%%%%%%%%%%%%%%
%%%%%%%%%%%%%%%%%%%%%%%%%%%%%%%%%%%%%%%%%%%%%%%%%%%%%%%%%%%%
%%%      											    Title Page Information							                      %%%
%%%%%%%%%%%%%%%%%%%%%%%%%%%%%%%%%%%%%%%%%%%%%%%%%%%%%%%%%%%%
%%%%%%%%%%%%%%%%%%%%%%%%%%%%%%%%%%%%%%%%%%%%%%%%%%%%%%%%%%%%
\title{\vspace{-2.4cm}ECON 311: Intermediate Macroeconomic Theory}
\author{University of Northern British Columbia}
\date{WINTER 2026}
%%%%%%%%%%%%%%%%%%%%%%%%%%%%%%%%%%%%%%%%%%%%%%%%%%%%%%%%%%%%
%%%%%%%%%%%%%%%%%%%%%%%%%%%%%%%%%%%%%%%%%%%%%%%%%%%%%%%%%%%%
%%%         								  Double and Single Space Commands                                                            %%%
%%%%%%%%%%%%%%%%%%%%%%%%%%%%%%%%%%%%%%%%%%%%%%%%%%%%%%%%%%%%
%%%%%%%%%%%%%%%%%%%%%%%%%%%%%%%%%%%%%%%%%%%%%%%%%%%%%%%%%%%%

\newlength{\defbaselineskip}
\setlength{\defbaselineskip}{\baselineskip}
\newcommand{\setlinespacing}[1]{\setlength{\baselineskip}{#1 \defbaselineskip}}
\newcommand{\doublespacing}{\setlength{\baselineskip}{1.3 \defbaselineskip}}
\newcommand{\singlespacing}{\setlength{\baselineskip}{1\defbaselineskip}}
\renewcommand{\thesubsection}{\arabic{subsection}}
\newtheoremstyle{dotlessS}{}{}{\itshape}{}{\color{dg}\bfseries}{}{ }{}
\theoremstyle{dotlessS}
\newtheorem{theorem}{Theorem}
\newtheorem{lemma}{Lemma}
\newtheorem{proposition}{Proposition}
\newtheorem{assumption}[theorem]{Assumption}
\newtheorem{corollary}{Corollary}
\newtheorem{definition}{Definition}
\newenvironment{example}[1][Example]{\begin{trivlist}
\item[\hskip \labelsep {\bfseries #1}]}{\end{trivlist}}
\newenvironment{remark}[1][Remark]{\begin{trivlist}
\item[\hskip \labelsep {\bfseries #1}]}{\end{trivlist}}

\newcommand\Pitem{%
  \addtocounter{enumi}{-1}%
  \renewcommand\theenumi{\arabic{\enumi}'}%
  \item%
  \renewcommand\theenumi{\arabic{enumi}}%
}
%%%%%%%%%%%%%%%%%%%%%%%%%%%%%%%%%%%%%%%%%%%%%%%%%%%%%%%%%%%%
%%%%%%%%%%%%%%%%%%%%%%%%%%%%%%%%%%%%%%%%%%%%%%%%%%%%%%%%%%%%
%%%           								                    Set Margins Here                                                                  %%%
%%%%%%%%%%%%%%%%%%%%%%%%%%%%%%%%%%%%%%%%%%%%%%%%%%%%%%%%%%%%
%%%%%%%%%%%%%%%%%%%%%%%%%%%%%%%%%%%%%%%%%%%%%%%%%%%%%%%%%%%%
\textwidth=6.5in
\textheight=9in
\oddsidemargin=0.10in
\evensidemargin=0.10in
\topmargin=-0.5in
%%%%%%%%%%%%%%%%%%%%%%%%%%%%%%%%%%%%%%%%%%%%%%%%%%%%%%%%%%%%%
%%%%%%%%%%%%%%%%%%%%%%%%%%%%%%%%%%%%%%%%%%%%%%%%%%%%%%%%%%%%%
%%%          									  Actual Document Begins Here                                                               %%%
%%%%%%%%%%%%%%%%%%%%%%%%%%%%%%%%%%%%%%%%%%%%%%%%%%%%%%%%%%%%%
%%%%%%%%%%%%%%%%%%%%%%%%%%%%%%%%%%%%%%%%%%%%%%%%%%%%%%%%%%%%%
\begin{document}
\maketitle
\vspace{-1cm}
\noindent\makebox[\linewidth]{\rule{\textwidth}{1.5pt}}

\section{Contact Information}
\subsection*{Instructor: Leandro Freylejer}
\textbf{Office:} Teaching and Learning Building, Room 10-4536\\
\textbf{E-mail:} leandro.freylejer@unbc.ca, \textbf{Phone:} 250-960-5302 \\
\textbf{Office Hours:} Tuesdays 11:45am-12:45pm, Thursdays 1:00pm-2:00pm
\section{Course Information}
\textbf{Room:} Agora 7-158, \textbf{Time:} Tuesdays, Thursdays 8:30 am-9:50 am \\
\textbf{Course website:} https://moodle.unbc.ca/


\section{Overview}
This course provides an introduction to modern macroeconomic theory at an intermediate level. Topics covered in this class include: measurement of aggregate economic activity; long-run economic growth; the quantity of money and its effect on inflation; as well as short-run macroeconomic analysis, including stabilization policies. The focus of this course will be mainly theoretical. However, I will try to connect the theory with recent economic issues whenever possible by assigning additional readings from the ``popular'' press and/or including empirical applications in assignments.    

\section{Textbook}
The required textbook for this course is \textbf{\textit{Macroeconomics}, by Charles I. Jones}, Sixth Edition (Fifth and Fourth Editions are also acceptable).  
\section{Lectures, Tutorials and Office Hours}

\subsection*{Lectures}
Lectures will take place on Tuesdays and Thursdays from 8:30 am to 9:50 am in Room 7-158. In most weeks, I will post lecture slides on the course website at least one and a half hours before class. If necessary, corrected slides will be posted as soon as possible. Slides are not substitutes but complements to class attendance. I also encourage students to complete readings before class. This facilitates in-class discussion and participation. I encourage you to use laptops or tablets only for the purposes of taking notes. I ask you to please turn off your cell phones and other communication devices before entering the classroom. Please be mindful of other students around you. Any conduct that is deemed to be disruptive to a learning environment will not be tolerated.  \\ 

\subsection*{Tutorials}
There will be two tutorials in this class scheduled for February 12th and April 9th during class time. In these tutorials, I will go over solutions to assignments and additional practice questions, which will posted on Moodle prior to the tutorial. It is important that you attend tutorials and come prepared to ask questions. Tutorials are scheduled for the class immediately preceding exams, and therefore, they might be your last chance to ask questions before being tested on the material. 

\subsection*{Office Hours}
I will hold office hours on Tuesdays 11:45am-12:45pm and Thursdays 1:00pm-2:00pm. A lot of the later material in this course builds on earlier material. Therefore, it is very important that you do not let major gaps in your understanding drag on. You are strongly encouraged to come to office hours as early as possible and as \textbf{often} as you need.\\
\\
\noindent If for some specific reason you cannot attend office hours, I am available by appointment only. Please do not drop-in unannounced.


\section{Evaluation}
\subsection*{Due Dates and Final Grade Breakdown}
The final grade for the course consists of \textbf{term work worth $60\%$} and a \textbf{comprehensive final exam worth $40\%$}.
\\
\begin{table}[!h]
\begin{minipage}[h]{1\textwidth}
\centering
\resizebox{10cm}{!}{%
\begin{tabular}{cccc}   
\hline
					     & Value                & Due date   \\  
\hline
\hline
First Assignment           & $10\%$           & February 10th, in class   \\
Midterm Exam              & $25\%$           & In-class: February 24th                       \\ 
Critical Review             & $15\% $          & In-class: March 26th                 \\
Second Assignment      & $10\%$           & April 7th, in class                  \\
Final Exam                   & $40\%$           & TBD                                                 \\
\hline
\end{tabular}}
\end{minipage}
\end{table}	

\noindent Throughout the semester, I might assign additional problems that I will solve in class. These problems will not be graded but I expect that you attempt to solve them before class.\\ 
\subsection*{Assignments}
There will be two graded assignments in this course. In these assignments, I will ask you to solve a combination of theoretical and empirical problems. Theoretical problems will be similar to the types of problems that will appear on exams. In the empirical problems, I will ask you to analyse macroeconomic data through the lens of the theory studied in class. I will post these assignments on the course website \textit{at least} two weeks before the due date. I encourage you to start working on them as soon as they are posted. 

\subsection*{Critical Review}
To be successful in the workforce, you will need to be able to critically evaluate technical arguments and convey your opinion in writing. The aim of this assignment is to provide you with these skills. For the critical review, I will assign a particular topic relating class material to a contemporaneous economic issue (e.g. Covid and the economy, tariffs, etc.). Students are responsible to do their own independent research on this topic in advance. On the date of the critical review, students will bring a report on their research which they will use to critically analyse an article on the assigned topic. In addition to assigning the topic in advance, I will provide students with a detailed grading rubric, which will be used to evaluate the submission, before the date of the critical review. However, I will not circulate the article. The critical review must be completed in class using pen and paper. I will provide more information about the critical review closer to the due date.      
    
\subsection*{Exams}
There will be two exams in this class. A midterm exam scheduled for February 24th during class time, and a \textbf{comprehensive} 3-hour final exam to be scheduled by the University during the examination period. In these exams, I will ask you to solve a combination of short/long-answer questions. I will provide you with more information regarding the format of these exams closer to the examination dates.  

\section{Administrative Details}
\subsection*{Course web-page and e-mail policy}
Students should check Moodle (moodle.unbc.ca) frequently for announcements, assignments, additional resources and other information for this course.  Moreover, I often post additional practice material, such as past exams with answer keys, before midterm and final exams.\\
\\   
I will do my best to respond to e-mails within 24 hours on a weekday, 48 hours on a weekend. Please use official ``@unbc.ca" email accounts for all communication. I only respond to e-mails posing questions that can be answered in no more than three sentences. For more involved questions, I will ask you to attend my office hours or ask questions during lectures. Also, I do not reply to e-mails that request information that can be found in the website or the syllabus.

\subsection*{Academic Misconduct}
I do not tolerate any form of academic misconduct, including \textit{plagiarism}. Any student caught engaging in such activities will be subject to academic discipline ranging from a mark of zero on the test or assignment to dismissal from the university as outlined in the academic handbook.

\subsection*{Late Submissions}
Late assignments and critical reviews are subjected to a penalty of $8$ percentage points per day, and they will not be accepted two days after the deadline has passed (including weekends). If you cannot submit an assignment on time, for an exceptional reason out of your control, you should email me as soon as you are able to. We will discuss your particular situation, and hopefully, come up with a mutually satisfactory solution, which is also fair to other students in the class who submitted their assignment on time. When proper documentation can be obtained, you should be prepared to provide such documentation.    


\subsection*{Missing an Exam}
A missed midterm or final exams will receive a zero grade, unless in exceptional circumstances when for documented serious reasons a student is not able to write an exam. In this case, I will shift the weight of the midterm towards the final exam or allow a student to write a make-up final exam. If you miss an exam, you should email me as soon as you are able to. We will discuss your particular situation, and hopefully, come up with mutually satisfactory solution, which is also fair to other students in the class who did not missed the exam. In the case when proper documentation can be obtained, you should be prepared to provide such documentation. 

\subsection*{Grade Appeals}
Appeals regarding the grading of assignments, critical review or exams must be submitted to me by email within two weeks of your receipt of the graded work. You must clearly and in detail explain the reason for your appeal. I will re-grade the entire work. Note that this may lead to a higher or lower overall grade. This policy does not apply to situations where there is a clear mistake on our part (e.g. I added points incorrectly)

\newpage

\section{Tentative schedule of Topics and required readings}
\begin{table}[!h]
\begin{minipage}[h]{1\textwidth}
\centering
\resizebox{18cm}{!}{%
\begin{tabular}{|c|c|c|c|}
\hline  
Class                                                      & Topic                                             & Chapters/Readings   & Description \\
\hline
					                             &									   &                                              &                                                          \\
\multirow{ 2}{*}{1-2}                          &\multirow{ 2}{*}{Introduction}        &\multirow{ 2}{*}{1-2}         & Introduction to macroeconomics and   \\
	                                                      &								         &                                            &   measurement of aggregate economic output	   \\[1.4ex]
\hline                 																												
			                      			&						                                &                                              &                                                          \\
3-5                                          & \multirow{ 5}{*}{Long-Run}                 & 3-4                                       & Introduction to long-run macroeconomics       \\[1.4ex]
6-8                                          &                                                         &5                                            & The Solow Growth Model                \\[3ex]
\multirow{ 2}{*}{9-10}               &                                                         &6,                                              &\multirow{ 2}{*}{ Long-run growth}  \\
                                                        &                                                         & GEHT (146-152, 156-166)         &                                                                      \\[1.4ex]
\hline
						&									&                                                  &                                                          \\
11                                 & Tutorial I                                           &                                                 & Solution to Assignment 1 and Practice Questions       \\[1.4ex]
\hline
						&									&                                                  &                                                          \\
12                             & Midterm Exam                                     &                                    &                  \\[1.4ex]
\hline
			                      			&						                                &                                              &                                                          \\
13 (1st part)                         & Long-Run                                                       & 8 (pp.216-230-- 5Ed., pp.231-248-- 6Ed.)                          &Inflation in the long-run     \\[1.4ex]						\hline
			                      			&						                                &                                              &                                                          \\
13(2nd part)-14                           &   \multirow{ 6}{*}{Short-Run}          & 9                                           & Introduction to short-run macroeconomics \\[3ex]
\multirow{ 2}{*}{15-18}				&									       &\multirow{ 2}{*}{11, 12}                                    & The IS-MP framework, fiscal policy,  \\   
                                                        &                                                              &                                                                            & and the Phillips Curve \\  [3ex]        
19-21      				&									&13                                         &The AS-AD framework and monetary policy            \\[3ex]
\multirow{ 2}{*}{22-23}                                 &                  &   10, 14                                            & Recent Economic Crisis:                                   \\
                                                                  &                  &  RDE (pp.165-172), TBA                                           &  Great Recession (2008), Covid Crisis                                                     \\[1.4ex]
\hline
						&									&                                                  &                                                          \\
24                                &   Tutorial II                                       &                                              & Solution to Assignment 2 and Practice Questions                    \\[1.4ex]
\hline
\end{tabular}}
\end{minipage}\captionsetup{font={normal}}
\captionsetup{width=0.99\textwidth}  
\caption*{Note: additional readings will be posted in the Additional Material folder in Moodle. GEHT refers to ``Good Economics for Hard Times'', by Banarjee and Duflo. RDE refers to ``The Return of Depression Economics'', by Paul Kurgman}
\end{table}	


\newpage

\section{Additional University Resources}

The Academic Success Centre provides students with FREE access to academic support services: 
\begin{itemize}
\item Tutoring (by appointment, asynchronous online, or drop-in)
\item Personalized study skills assessments 
\item Peer-led course supports 
\item Downloadable handouts 
\item Access to self-assessment sites 
\item Customized programs and workshops
\end{itemize}
ASC services are available in person at the Prince George Campus and online for distance, regional, or physically distanced students. \\
\\
For more information see the Academic Success Centre website at \url{www.unbc.ca/asc}, to book an appointment visit them on the first floor of the library, email asc@unbc.ca, or phone 250-960-6367. Toll-free: 1-888-440-3440 (wait on the line and ask for the ASC).



\subsection*{Access Resource Centre}
The Access Resource Centre (ARC) provides service to students with documented disabilities or health conditions, ranging from permanent to temporary, including but not limited to chronic health issues, hearing and visual impairments, learning disabilities and attention deficit/hyperactivity disorders, mental health and neurological disabilities, and mobility and other physical disabilities. ARC staff are available by appointment to assess specific needs, provide referrals, and arrange appropriate academic accommodations to assist students in achieving their academic goals. Students who may have a need for academic accommodation are encouraged to contact ARC by email at arc@unbc.ca, by phone at 250-960-5682 (toll free 1-888-960-5682), or in person at 5-157. 


\subsection*{Health and Wellness}
Everyone can experience mental health concerns or stressful events that can interfere with our academic performance and negatively impact our daily activities.
All of us can benefit from support during times of struggle.  If you or anyone you know experiences academic stress, difficult life events or feelings of anxiety, depression or suicidal thoughts, the UNBC Wellness Centre team is here to help. Their services are free for UNBC Students and appointments are available. You can learn more about confidential mental health services available on and off campus at: https://www2.unbc.ca/counselling  OR https://www2.unbc.ca/medical-clinic. Remember that getting help is a smart and courageous thing to do - for yourself, for those you care about, and for those who care about you. Asking for support sooner rather than later is almost always helpful. 




\end{document}

